\section{Introduction}
A maintainer receives a patch for a price-history error and reruns the test from the repair ticket. The test passes. Another caller can read the same pool's history through a different contract interface. Should the maintainer accept the patch for that caller too? This was a concrete distinction in an earlier study of Ekubo. Some patches returned the expected history through the library interface while the direct contract interface still returned an incorrect answer. One passing test settled the question it asked. The maintainer's decision depended on a second question.

Solidity makes the distinction observable because different callers can depend on the same persistent contract state. A change that prevents a reported failure may also interrupt a legitimate operation. Conversely, a patch can keep ordinary workflows running while leaving the audited defect visible under another use. EVMPatch made both aspects measurable by checking the effect of Ethereum patches against harmful and legitimate historical transactions~\cite{rodler2021evmpatch}. For a repair agent that reads a regression example in the ticket, we want to distinguish learning what the contract should do from fitting what that example happens to exercise.

Development records give reasons to expect both incomplete repairs and ceilings. On one finding, a second check rejected patches that had passed the supplied example. On another, new checks rejected patches that had already passed a battery of checks and the ordinary-use suite. A different configuration on the latter finding passed even the expanded battery across all its counted sessions. These observations establish that the choice of measurement matters on these cases. They leave open whether changing the example given to a repair agent would change the patch.

The missing comparison is specific. For the same audited defect, keep the ticket, code, working environment and grader fixed. Vary only which of two legitimate examples a maintainer includes. Under either assignment, ask separately whether each behaviour now works. A finding can contribute to a benchmark of reported repairs even when its public history supports only one observable obligation. Counting a second copy of that same obligation as evidence of transfer would make the apparent breadth of a repair depend on how many equivalent tests a curator wrote.

We therefore ask three questions with explicit denominators. \textbf{RQ1 (reported behaviour).} On admitted findings, how often does a repair configuration deliver a \emph{working repair of the reported behaviour}? The submitted patch builds, preserves the specified ordinary use, and corrects the defect shown in example A for a reason supported by the recorded assessment. \textbf{RQ2 (another use).} Among all assigned A attempts on findings with a qualified hidden behaviour H, what fraction satisfies both A and H while preserving ordinary use, and how much lower is that fraction than the primary rate on those same attempts? \textbf{RQ3 (example choice).} Among findings with two legitimate alternatives A and B and a common withheld H, does randomly showing A or B preferentially improve the behaviour illustrated? RQ1 and RQ2 use standard reports showing A. RQ3 adds the alternative assignment.

This manuscript presents two contributions in hand and a prospective experiment:
\begin{itemize}
\item A comparison of the evidence behind three archived development results, including their different denominators and when the graders were constructed.
\item A case-admission and measurement protocol, with an account of which public findings have reached each qualification step and which evidence is still missing.
\item A registered primary outcome and a design for comparing repair configurations and the identity of the supplied example. The performance estimates are due after admission, assignment and grading.
\end{itemize}

\section{What earlier repair studies measure}
\paragraph{Changing the tests available during repair.}
Smith et al.\ evaluated patch overfitting using held-out test suites~\cite{smith2015overfitting}. Motwani et al.\ came closer to our intervention: within defects, they sampled failing tests while holding passing tests fixed and examined patch quality as the number of failing tests changed~\cite{motwani2022quality}. Modern agents also use tests as feedback. Ahmed et al. studied test overfitting on SWE-bench, including access to hidden developer tests~\cite{ahmed2026overfitting}. TDFlow contrasts developer-supplied tests with tests produced by the system~\cite{han2026tdflow}. These studies motivate controlling the amount and origin of test information. We fix the number of reported examples at one and estimate two directional contrasts within each defect. This lets an intrinsically easier behaviour succeed often under both assignments without being mistaken for an advantage from showing it.

\paragraph{Measuring a security repair.}
Vul4J links real Java vulnerabilities to historical revisions and executable evidence~\cite{bui2022vul4j}. SEC-bench evaluates end-to-end agents on real security tasks~\cite{lee2025secbench}. EVMbench already evaluates smart-contract patching using existing project tests and hidden checks~\cite{wang2026evmbench}. The PatchBench preprint examines how security and semantic checks change patch assessments in C and C++~\cite{shen2026patchbench}. Independent testing, Solidity patching, and failure under stronger grading thus have direct precedents. The proposed contribution concerns \emph{which reported behaviour to illustrate} and the conditions under which a second observation can support an inference about another use. We limit this distinction to the studies compared here.
